\documentclass[11pt]{article}
\usepackage{ideastyle}

\begin{document}

\ideatitle{Wanderlings}{Pikmin Bloom, but your little companions are AI agents that scout the real world for you.}

\section{The Idea}
In Pikmin Bloom, you walk, flowers bloom along your path, and you send little
creatures on expeditions. Wanderlings keeps that cozy walking loop but makes the
creatures \emph{real agents}: each one can plan, use tools, and come back with
something useful.

Send your Scout Wanderling to that park you've never visited. It ``travels'' there
(taking real time based on distance), looks the place up using real map and
Wikipedia data, decides if you'd like it, and returns with an illustrated postcard
(Grok Imagine) and a short voice memo (Grok Voice): \emph{``Found a waterfall
12 minutes away. Mostly flat path. Want me to lead you there?''} Then it guides you
there while your walk blooms on the map.

\textbf{Navigation angle:} instead of the shortest path, you navigate by
curiosity. Your agents discover places and routes worth walking.

\section{The Agents}
\begin{tabular}{@{}P{0.2\textwidth}P{0.74\textwidth}@{}}
\toprule
\textbf{Scout} & Finds nearby places you've never been and reports back. \\
\textbf{Storyteller} & Tells you the history of landmarks you pass (from Wikipedia, not made up). \\
\textbf{Pathfinder} & Plans a scenic or step-free route to wherever the Scout found. \\
\textbf{Forager} & Finds coffee, food, water fountains, and benches along your walk. \\
\bottomrule
\end{tabular}

Each agent levels up and changes appearance based on where you walk
(parks make it leafy, libraries make it bookish), like Pikmin's decor variants.

\section{Track Fit}
\begin{tabular}{@{}P{0.28\textwidth}P{0.66\textwidth}@{}}
\toprule
\req{Built with Cursor}{Agent logic, map, and game UI built in Cursor.}
\req{Grok Voice / Imagine}{\textbf{Imagine} for postcards and creature art; \textbf{Voice} for agent voice memos. Uses both.}
\req{Space data}{Not used.}
\req{Navigation theme}{Discovery-driven navigation; agents as scouts.}
\req{Also eligible}{Big Red, Software, Design, People's Choice. \textbf{Photon} if agents text you in iMessage.}
\bottomrule
\end{tabular}

\section{Features}
\subsection{MVP}
\begin{itemize}
  \item Map with your walked path drawn as a trail of flowers.
  \item Three agents (Scout, Storyteller, Pathfinder), each with its own personality.
  \item \textbf{Expedition:} pick a nearby place; the agent researches it with real tools (map data, Wikipedia) and returns after a short timer.
  \item Return gift: Grok Imagine postcard + Grok Voice memo + ``Take me there'' button.
  \item Walking directions to the place, with the trail blooming as you go.
\end{itemize}
\subsection{Stretch}
\begin{itemize}
  \item Agents evolve based on the kinds of places you visit.
  \item Daily quests the agents invent (\emph{``Visit 2 places you've never been''}).
  \item Agents text you in iMessage with Photon (\emph{``I'm back from the gorge!''}).
  \item Friends' agents meet when you walk near each other.
\end{itemize}

\section{Tech Stack and Data}
\begin{itemize}
  \item Frontend: mobile web app (PWA), Leaflet map, browser geolocation.
  \item Agents: an LLM with tool calls (Grok chat API fits the track), tools for:
    nearby places (OpenStreetMap Overpass API), place summaries (Wikipedia API),
    walking routes (OSRM / Google Directions).
  \item Media: Grok Imagine (postcards, creature art), Grok Voice (memos).
  \item Backend: small Node or Python server holding API keys and expedition timers.
  \item Optional: Photon Spectrum for iMessage.
\end{itemize}

\section{Demo Plan (2 minutes)}
\begin{enumerate}
  \item Show Saturday's real walk around campus, blooming on the map.
  \item Send the Scout on an expedition live (shortened timer for the demo).
  \item It returns: play the voice memo and show the postcard.
  \item Tap ``Take me there''; the Pathfinder plans the route.
\end{enumerate}

\section{Limitations and Risks}
\begin{itemize}
\limitation{Indoor judging, no walking}{You can't go on a real walk at the judging table.}{Record real walks Saturday and add a ``replay walk'' mode; expeditions themselves work anywhere.}
\limitation{Background location on the web}{Web apps can't reliably track location with the screen off, especially on iPhone. Real walking games are native apps for this reason.}{Keep the app open during walks for the MVP; say a native app is the next step.}
\limitation{Agents making things up}{An LLM may invent facts about a place (``famous waterfall'' that doesn't exist).}{Agents may only report facts returned by their tools; show the sources on the postcard.}
\limitation{Postcards aren't real photos}{Grok Imagine illustrations won't look exactly like the place.}{Keep a clear illustrated, storybook style so nobody mistakes them for photos.}
\limitation{Copyright and trademarks}{Pikmin is Nintendo/Niantic's. Using its name, art, or creatures is not allowed.}{Original creatures and name; say ``inspired by walking games'' in the pitch.}
\limitation{Scope creep}{Games invite endless features (levels, items, friends).}{Lock the MVP to one loop: walk, send, return, go. Everything else is stretch.}
\limitation{Cost and speed}{Each expedition calls an LLM several times plus image and voice generation.}{Cache results per place; pre-generate the demo expedition as a backup.}
\limitation{Location privacy}{The app stores where people walk.}{Keep walk history on the device; send the server only the place being researched.}
\end{itemize}

\section{Scorecard}
\begin{tabular}{@{}P{0.25\textwidth}P{0.08\textwidth}P{0.6\textwidth}@{}}
\toprule
\rating{Demo-able}{4}{Cute and memorable; walking must be replayed.}
\rating{Buildable in 24h}{3}{Several APIs, but each piece is standard.}
\rating{Navigation theme}{4}{Navigation by curiosity; slightly playful take.}
\rating{Wow factor}{5}{Charming creatures that really do things.}
\rating{Technical risk (5 = riskiest)}{3}{Moderate (many integrations).}
\bottomrule
\end{tabular}

\section{Verdict}
\textbf{Strong contender, especially for People's Choice and Design.} It uses
both Grok APIs naturally, and the agents do real work instead of only chatting.
Keep the MVP to a single expedition loop.

\end{document}
